\section{A compact positive-domain companion}
\label{sec:compact-center}

This companion has ten strictly positive integer variables.  Its domain
must be kept separate from the unrestricted-integer equations in the main
construction.  It combines the unsquared strong trace of the pinned
\texttt{163-trace-radical} contribution with the inverse-Lucas correction
of \texttt{164-unsquared-factors-linear-selectors}~\cite{trace163,inverse164}.  A new square transport
weight then permits a polynomial root center and lowers the exact degree.

\begin{center}
\begin{tabular}{lrrr}
\hline
Circuit & Variables & Degree & Binary operations\\
\hline
Trace parent & 10 & 232 & 87\\
Inverse correction with trace & 10 & 232 & 85\\
Square transport with trace & 10 & 232 & 85\\
Centered companion & 10 & 224 & 90\\
\hline
\end{tabular}
\end{center}

The degree 224 point is a measured degree/cost tradeoff.  The claimed result
is neither a global minimum nor an assertion of literature-wide priority.
This compact branch retains its inherited native/K215 calibration.  Its
ordinary proof is separate from its finite algebraic checker.

\subsection{Entire polynomial}

The independent variables are
\[
J,F,\alpha,i,T,w,s,n,\eta,\zeta\in\mathbb Z_{>0}.
\]
The four fixed coefficients \(\beta,K_0,M_C,M_F\) are the whole
indexed BFS943/radix67 tuple, with \(\beta=2^{590069}-1\).  Their complete
finite recipe, 436-character seed, numerical payloads, and independent
reconstruction program are included.  All raw and compressed hashes agree
with the pinned parent.  The largest port has 592,884 bits; the four raw
payloads occupy 295,388 bytes.  Every symbol below is a polynomial
abbreviation, not another quantified variable:
\begin{gather*}
h=\beta J,\quad q=h+1,\quad C_0=q-F-\alpha,\\
X=qw,\quad Y=qs,\quad k=\eta+\zeta,\quad c=kY+\eta,\\
A=Y(X+1)+1,\quad \Delta=A^2-1,\quad H=4A-5,\\
v=XY^2,\quad \ell=vk=(XY)(kY),\quad E=1+\ell(\ell+k),\\
R=[h(q-F)+\alpha](q^2-1)+(M_C+qM_F)J,\\
B=(K_0+w)C_0-F,\quad U=ic^2,\quad a=\Delta U^2E,\\
F_0=1+a,\quad f=1+2a,\quad G=EF_0,\quad Q=4\Delta aF_0,\\
b=cT-R,\quad V=fb-c,\quad K=X(3c-2X)+2,\\
e_6=6(k-R-1)-\ell(k^2-4),\quad D=6q^6,\quad M=hDH.
\end{gather*}
For a weight \(L\), set
\[
Z_0=hD\bigl[H(n-LB)+K\bigr]-H\bigl[DB+h(e_6-6q^4s)\bigr].
\]
The centered companion uses
\[
L=b^2,\qquad C=AUE,\qquad S=CV,\qquad Z=Z_0-MS
\]
and the single equation
\begin{equation}
\boxed{\ P=Q\bigl(Z^2-M^2GV^2\bigr)+(M^2G-Z^2)=0.\ }
\label{eq:compact-centered}
\end{equation}
The two 85-operation controls use \(Z=Z_0\) and respectively
\(L=bV\) or \(L=b^2\).  All three literal DAGs are supplied.  Fixed
powers are multiplication gates; no variable exponent, recurrence, floor,
or unspecified coefficient occurs in the polynomial.

For proof notation only, put
\[
e_*=\frac{e_6}{6},\quad \rho=\frac{e_*}{q^6},\quad
\sigma=\frac{s}{q^2},\quad t=\frac{B}{h},\quad W=1+hL,
\quad\kappa=\frac{K}{H},\quad z=\frac{Z}{M}.
\]
Without yet asserting integrality, the centered carrier identity is
\begin{equation}
z=n-S+\kappa+\sigma-Wt-\rho.
\label{eq:compact-carrier}
\end{equation}
For an uncentered control, omit \(S\).

\subsection{The square transport theorem}

The root extraction, rational adjacent-Pell descent, signed folding, and
third-index exclusions in the trace source do not use the old cubic, the
weight, or the carrier center.  Before assuming any sign for
\(B,C_0,R,b,V\), their outputs are
\begin{gather*}
G=g^2,\quad E=\tau^2,\quad F_0=f_0^2,\quad
z=\varepsilon gy,\quad QV^2-(Q-1)y^2=1,\\
c=\psi_A(p),\quad p\ge5\text{ odd},\quad
k=2\psi_{1+2v}(m),\quad m\ge2,\\
c>8XY^3\ge8q^4,\quad c>X,\quad K>0,\quad
b\ne0,\quad \operatorname{sgn}V=\operatorname{sgn}b,\\
Q>f^2,\quad c<f/2,\quad |V|>8Q^2,\quad
|b|>8f^3-\tfrac12.
\end{gather*}
Here \(g,\tau,f_0,y\) are positive integers and
\(\varepsilon\in\{1,-1\}\).  The sequence \(\psi_A\) has initial
values 0,1 and recurrence \(\psi_A(j+2)=2A\psi_A(j+1)-\psi_A(j)\).
If the auxiliary index were 3, its two sign cases would give, respectively,
\[
f(b-4\Delta f)=c-(4\Delta+3),\qquad
f(b+4\Delta f)=c+(4\Delta+3):
\]
Both right sides lie strictly between 0 and \(2c<f\), whereas each
left side is a multiple of $f$, a contradiction.

We now establish the new bound needed for the center.  The identities give
\[
\frac{a^2}{g^2}=\Delta U^2\frac{a}{1+a}>9,
\qquad \frac{C^2}{G}<\frac{A^2}{A^2-1}<4.
\]
Thus \(g<f/6\) and \(C<2g\).  Also
\(y<\tfrac32|V|\) and \(|V|<\tfrac32f|b|\).  For either root sign,
\begin{align*}
|\varepsilon gy+CV|
&\le gy+C|V|<\tfrac72g|V|
 <\tfrac78 f^2|b|\\
&<\tfrac7{32}b^2<\tfrac14b^2.
\end{align*}
We used \(|b|>4f^2\), a consequence of the stronger bound above.

The inverse-Lucas correction is nonnegative before rank sign recovery:
\(vk(k^2-4)/6\ge0\), because \(k\ge2\).  Consequently
\[
e_*\le k-R-1<b,\qquad \rho<\frac{|b|}{q^6}<\tfrac14b^2.
\]
If \(B\le-1\), then \(Wt<-L\).  Solving
\eqref{eq:compact-carrier} for positive \(n\) gives
\[
n<\tfrac14L-L+\tfrac14L<0,
\]
a contradiction.  Therefore \(B\ge0\), and then \(C_0>0\).
This treats all integer ranks \(R\), including negative and zero ranks.
The missing adverse-cubic estimate identified in the trace source's open
\(b^2\) question is unnecessary for the inverse correction.

\subsection{Soundness and the prime 3 ordering}

After recovering \(C_0>0\), the calibrated mask bound gives
\(3q<R<q^4-2\), hence \(c>2R\) and \(b,V>0\).  Since
\[
f-1=2\Delta i^2c^4E,\qquad V\equiv-c\pmod f,
\qquad V\equiv-R\pmod c,
\]
the same normalized-rank theorem gives \(p=R\).  In particular \(R\)
is positive odd, \(q\) is even, and \(J\) is odd.  The new center
\(S\) is an integer polynomial, so integrality of \(z\) implies
integrality of \(z+S\).  The remaining carrier argument is therefore
the one in the inverse-Lucas source; the weight enters only through
\(W\equiv1\pmod h\).

For clarity, its essential arithmetic is the recurrence expansion
\[
k=2m+8vC_m+32v^2Z_m,
\qquad C_m=\binom{m+1}{3},\qquad Z_m\in\mathbb Z_{\ge0}.
\]
Writing \(d=2m-R-1\) and \(A_m=C_m+4vZ_m\), one obtains
\[
3(e_*-d)=4v^2\bigl[24Z_m-A_m\{k^2+2mk+(2m)^2-4\}\bigr].
\]
The integer $k$ is even, so $k(k^2-4)/6$ and hence $e_*$ are integers.
Carrier integrality gives \(q^6\mid e_*-q^4s\), hence
\(q^4\mid e_*\).  Prime-by-prime valuations of the last identity give
\(q^4\mid e_*-d\), including the possible loss of one factor 3.
The bound \(1\le m<R<q^4\) forces \(d=0\).

One must not yet conclude \(q^6\mid e_*\).  First the whole-product
projection from \(H\mid6hK\), with its retained strict size bound,
gives \(X=2^R\).  Since \(q\mid X\), \(q\) is a power of 2.  Only
then can the factor 3 be canceled and \(q^6\mid e_*\) recovered.
The remaining congruences give \(q^2\mid s\), \(h\mid B\), and
integral positive \(\kappa\).  Thus
\[
X=2^R,\quad q^3\mid Y,\quad t=B/h\in\mathbb Z_{\ge0},\quad
(K_0+w)C_0-F=t(q-1).
\]
The inherited whole-tuple native/refutation calibration now yields a
finite syntactic ZFC refutation.  This step retains the entire calibrated
tuple, not merely its elementary coefficient inequalities.

\subsection{A cofinal positive lift of the centered equation}

Conversely a finite refutation supplies the canonical outer witnesses in
the two pinned source theorems.  In particular \(q\) is a power of 2,
\(R\equiv3\pmod4\), \(X=2^R\), \(q^3\mid Y\), \(C_0>0\),
and \(t\ge0\), while \(\rho,\sigma,\kappa\) are integers.
Their exact Pell-ratio interval supplies positive \(\eta,\zeta\).
Choose the strong seed with \(c^2\tau\mid I_0\), so
\(i=I_0/(c^2\tau)>0\), \(g=\tau f_0\), and \(f=2f_0^2-1\).

The retained adjacent-Pell completeness theorem supplies positive
\(V_*,y\) such that \(V_*+c+Rf^2\) is divisible by \(cf^2\).
The unsquared-trace coordinate
\[
T=\frac{V_*+c+Rf}{cf}
\]
is a positive integer, because \(c\mid f-1\).  For an uncentered
control take \(n=gy+Wt+\rho-\kappa-\sigma\).  The inherited bounds
\(|\rho|<g/3\) and \(\kappa+\sigma<g/2\) prove positivity.

For the centered equation the negative root is useful.  Put
\(u_0=U^2E\ge2\).  The exact identity
\[
(Q-1)C^2-QG=E\bigl[Q(u_0-1)-A^2u_0\bigr]
\]
is positive: \(Q>4\Delta^2u_0\) and \(4\Delta^2>A^2\).
Therefore \(C>g\sqrt{Q/(Q-1)}\).  If
\(D_0=(Q-1)C^2-QG\), then \(D_0\) is a positive integer and
\[
CV_*-gy>\frac{D_0V_*}{2C(Q-1)}\ge\frac{V_*}{2C(Q-1)}.
\]
There are cofinally many permissible \(V_*\).  The matrix
\[
\begin{pmatrix}2Q-1&2(Q-1)\\2Q&2Q-1\end{pmatrix}
\]
preserves the adjacent Pell norm, has determinant 1, and strictly
increases positive pairs.  A power of it is the identity modulo
\(cf^2\); iterating that power preserves the required congruence and
makes \(V_*\) unbounded.  No additional polynomial variable is used.

Choose \(V_*>2C(Q-1)(\kappa+\sigma+|\rho|)\) and set
\[
n=CV_*-gy+(1+hb^2)t+\rho-\kappa-\sigma.
\]
Then \(n\) is a strictly positive integer, the centered carrier is
\(z=-gy\), and \eqref{eq:compact-centered} vanishes.  This proves
positive zero-existence equivalence; it does not mistake an arbitrary
integer shear for an automorphism of the positive domain.

\subsection{Exact degree and independent finite verification}

The cancellation is an identity over a free polynomial ring:
\begin{align*}
P={}&(Q-1)Z_0^2-2(Q-1)MCVZ_0\\
&+M^2V^2\bigl[QE(U^2E-1)-C^2\bigr]+M^2G.
\end{align*}
The degrees of \(Q,M,V,U,E,Z_0,C,G\) are respectively
\(80,11,40,7,14,21,25,50\).  Its unique degree 224 term is
\(QM^2V^2U^2E^2\); all other summands have strictly lower degree.
With \(q_0=\beta J\), \(k_0=\eta+\zeta\), and
\(B_0=k_0sT-q_0^2(q_0-F)\), the complete highest homogeneous part is
\[
[P]_{224}=9\,2^{10}q_0^{96}w^{24}s^{56}i^{10}k_0^{32}B_0^2.
\]
The coefficient of
\(J^{96}w^{24}s^{58}i^{10}\eta^{34}T^2\) is
\(9\,2^{10}\beta^{96}>0\), proving the exact degree on the authentic
coefficient slice.  The uncentered controls have complete top
\[
[P]_{232}=-9\,2^{10}q_0^{100}w^{26}s^{58}i^{10}k_0^{32}B_0^2.
\]

The checker imports no producer.  It verifies every actual DAG gate by
independently specified polynomial cuts, all live inputs and coefficient
ports, the exact cancellation identity, and the complete leading forms.
Both alternative operators at every gate are deliberately introduced and
rejected:180 controls for the centered circuit and170 for each control.
Normal and optimized Python runs pass.  This authenticates the finite
artifacts; the unbounded mathematical arguments above remain essential.
